
%\needspace{4\baselineskip}
\color{uniblue}{\subsection{Токовая защита нулевой последовательности (ТЗНП)}}
\color{black}

\begin{enumerate}[label=\arabic{section}.\arabic{subsection}.\arabic{enumi}, labelsep=4pt, leftmargin=0pt, itemindent=57pt, itemsep=0pt, parsep=5pt] % если что можно крутить itemsep
%%%%%%%%%%%%%%%%%%%%%%%%%%%%%%%%%%%%%%%%%%%%%%%%%%%%%%%%%%%  ОБЩИЕ СВЕДЕНИЯ %%%%%%%%%%%%%%%%%%%%%%%%%%%%%%%%%%%%%%%%%%%%%%%%

\item Общие сведения

\begin{enumerate}[label=\arabic{section}.\arabic{subsection}.\arabic{enumi}.\arabic*, labelsep=4pt, leftmargin=0em, itemindent=65pt, parsep=0pt]

\item
Токовая защита нулевой последовательности предназначена для селективной защиты от замыканий на землю в линиях среднего и высокого напряжения, а также для резервирования защит трансформаторов при замыканиях на землю в сети низкого напряжения, работающей с заземленной нейтралью.

\item
Устройство содержит четыре ступени ТЗНП с независимой выдержкой времени от КЗ на землю, реагирующие на первую гармонику тока нулевой последовательности.

\item
Основные характеристики функционального блока <<ТЗНП>>:
\begin{itemize}
    \item погрешность по углу при номинальных значениях тока и напряжения не более 5$^\circ$.
\end{itemize}

\item
Логическая часть реализует группу пусковых органов с независимыми уставками, предусматривает таймеры выдержек времени на срабатывание, логику пусков и блокировок, коммутацию внутренних пусковых сигналов для формирования выходного сигнала автоматического ускорения.

\item
ТЗНП может применяться в качестве:
\begin{itemize}
    \item селективной защиты от замыканий на землю в линиях среднего и высокого напряжения;
    \item резервной защиты трансформаторов при замыканиях на землю в сети низкого напряжения, работающей с заземленной нейтралью;
\end{itemize}

\item
В составе функционального блока <<ТЗНП>> предусмотрены следующие функции:
\begin{itemize}
\item четыре независимые направленные/ненаправленные ступени;
\item блокировка при броске намагничивающего тока;
\item орган направления мощности нулевой последовательности.
\end{itemize}

\item
Функциональный блок <<ТЗНП>> может быть полностью выведен из работы оперативно путем активации сигнала <<ОВ ТЗНП>>.

\end{enumerate}

%%%%%%%%%%%%%%%%%%%%%%%%%%%%%%%%%%%%%%%%%%%%%%%%%%%%%%%%%  СТУПЕНЬ ТЗНП %%%%%%%%%%%%%%%%%%%%%%%%%%%%%%%%%%%%%%%%%%%%%%%%

\item Ступени токовой защиты нулевой последовательности (ТЗНП 1 ст., ТЗНП 2 ст.)

\begin{enumerate}[label=\arabic{section}.\arabic{subsection}.\arabic{enumi}.\arabic*, labelsep=4pt, leftmargin=0em, itemindent=65pt, parsep=0pt]

\item
Ввод каждой ступени ТЗНП в работу на этапе параметрирования устройства осуществляется своим программным переключателем SGF1 <<Ввод\_функции>>. Информация о введенной ступени отображается выходным сигналом для каждой ступени (<<ТЗНП 1 ст: Ввод>>, <<ТЗНп 2 ст: Ввод>>, <<ТЗНП 3 ст: Ввод>>, <<ТЗНп 4 ст: Ввод>>).
\item
Все ступени <<ТЗНП>> могут быть выведены из работы оперативно путем активации сигнала оперативного вывода (<<ОВ ТЗНП>>), что характеризуется активным состоянием сигнала <<Оперативный вывод>> (<<ТЗНП 1 ст: Оперативный вывод>>, <<ТЗНП 2 ст: Оперативный вывод>>, <ТЗНП 3 ст: Оперативный вывод>>, <<ТЗНП 4 ст: Оперативный вывод>>) на выходе функции.
\item
При активном состоянии входного сигнала <<ТЗНП 1 ст на сигн>>, <<ТЗНП 2 ст на сигн>>, <<ТЗНП 3 ст на сигн>>, <<ТЗНП 4 ст на сигн>> действие соответствующей ступени ТЗНП переводится на сигнал.

\item
В качестве измерительных органов применяется ИО тока нулевой последовательности максимального действия, которые контролируют превышение тока НП заданной величины с коэффициентом возврата равным 0,95. Выбор режима работы ИО (независимая или обратнозависимая характеристика) осуществляется уставкой характеристики срабатывания <<Характеристика>>.
\item
При срабатывании ИО активируются соответствующие выходные сигналы :
\begin{itemize}
    \item <<ТЗНП 1 ст.: ИО 3I0>>;
    \item <<ТЗНП 2 ст.: ИО 3I0>>;
    \item <<ТЗНП 3 ст.: ИО 3I0>>;
    \item <<ТЗНП 4 ст.: ИО 3I0>>.
\end{itemize}
\item
При срабатывании ИО и выполнении условий логики работы схемы далее формируются выходные сигнала пуска:
\begin{itemize}
    \item <<ТЗНП 1 ст.: Пуск>>;
    \item <<ТЗНП 2 ст.: Пуск>>;
    \item <<ТЗНП 3 ст.: Пуск>>;
    \item <<ТЗНП 4 ст.: Пуск>>.
\end{itemize}
\item
При наличии пусковых условий функции запускается таймер Т1 зависимой/независимой выдержки времени на срабатывание. По истечении выдержки времени Т1, в течение которой величина тока постоянно превышала уровень тока срабатывания, активируются соответствующие выходные сигналы срабатывания:
\begin{itemize}
    \item <<ТЗНП 1 ст.: Срабатывание>>;
    \item <<ТЗНП 2 ст.: Срабатывание>>;
    \item <<ТЗНП 3 ст.: Срабатывание>>;
    \item <<ТЗНП 4 ст.: Срабатывание>>.
\end{itemize}
\item
При выборе режима с зависимыми характеристиками срабатывания активируется измерительный орган с зависимой выдержкой времени. Все характеристики (кроме характеристики RXIDG) в диапазоне значений тока от $ 1,0 \cdot I/I_{\textsl{ср}} $ до $ 1,1 \cdot I/I_{\textsl{ср}} $ имеют независимую выдержку времени на срабатывание ($t_{\textsl{ср.независ.}}= f (I = 1,1 \cdot I/I_{\textsl{ср}})$). При кратности тока выше 20 ($ I/I_{\textsl{ср}} \geq 20 $)  все зависимые характеристики переходят в независимые (кроме характеристики RXIDG) с выдержкой времени срабатывания, равной выдержки времени срабатывания при $ I/I_{\textsl{ср}} = 20 $ ($t_{\textsl{ср.независ.20}}= f (I = 20 \cdot I/I_{\textsl{ср}})$). C помощью временного коэффициента характеристики срабатывания ($K$) имеется возможность регулирования времени срабатывания времятоковых характеристик.

\item
Графики реализованных в устройстве обратнозависимых времятоковых характеристик приведены в приложении \ref{app:chars}.
\item
Ступень ТЗНП имеет возможность выбора направленности. Направленность ступени определяется уставкой SGF2 <<Направленность>>.  Ступень может работать как без контроля направленности, так и с контролем прямого либо обратного направления. Сигналы направления ступень получает от органа направления мощности <<ОНМ НП>>.

\item
Режим контроля направленности при неисправности цепей напряжения задается уставкой SGF4:
\begin{itemize}
    \item если уставка выставлена в положение <<Блокировка>>,  то при наличии активного сигнала неисправности ЦН осуществляется блокировка всех режимов направленности; 
    \item если уставка выставлена в положение <<Деблокировка>>,  то при наличии активного сигнала неисправности ЦН осуществляется разрешение работы во всех режимах направленности. 
\end{itemize}

\item
Ступень может работать с контролем напряжения НП, который осуществляется с помощью настройки уставки SGF6 <<Реж\_3U0>>. При активации контроля НП, пуск ступени возможен только в том случае, если утроенное напряжение нулевой последовательности превышает значение уставки и отсутствует сигнал неисправности от цепей напряжения, когда уставка SGF5 <<Реж\_БНН\_3U0>> находится в положении <<Блокировка>>.

\item
При срабатывании ИО напряжения НП активируются соответствующие выходные сигналы:
\begin{itemize}
    \item <<ТЗНП 1 ст.: ИО 3U0>>;
    \item <<ТЗНП 2 ст.: ИО 3U0>>;
    \item <<ТЗНП 3 ст.: ИО 3U0>>;
    \item <<ТЗНП 4 ст.: ИО 3U0>>.
\end{itemize}

\item
Работа ступени при неисправности цепей напряжения регулируется уставкой SGF5 <<Реж\_БНН\_3U0>>. Уставка может находиться в одном из двух состояний: <<блокировка>> или <<деблокировка>>. В режиме блокировки работа ступени запрещается до момента исчезновения неисправности БНН. В режиме деблокировки появление неисправности в цепях напряжения, контролируемых по входному сигналу <<БННш: Срабатывание>>, позволяет ступени функционировать без учета направления.

\item
При необходимости возможна блокировка ступени при броске тока намагничивания силового трансформатора от функции <<БНТ НП>>, это достигается путем перевода программного переключателя выбора режима контроля от БНТ SGF7 <<Реж\_БНТ>> в положение <<Предусмотрено>>.

\item
Вывод направленности ступени при автоматическом ускорении определяется уставкой SGF4 <<Ненапр\_АУ>>.

\item
Для блокировки чувствительных ступеней при работе в неполнофазном режиме в функции предусмотрена логика блокировки по дискретному входному сигналу <<Блок.чувств.ступ.>>. Блокировка
чувствительных ступеней формируется при введенном ключе <<Блок.чувств.ступ.>> в работу. Ступень блокируется, если уставка SGF8 <<Реж\_БЧС>> разрешает блокировку и на вход поступает сигнал <<Блок.чувств.ступ.>>.

\item
Алгоритм работы ступеней ТЗНП выполнен в соответствии с рисунком \ref{tznp:img1}. В таблице \ref{tznp:tbl1} приведены параметры, необходимые для настройки ступени ТЗНП.

\small
\begin{longtable}{|>{\centering\arraybackslash}m{5.3cm}|>{\centering\arraybackslash}m{3.3cm}|>{\centering\arraybackslash}m{4.2cm}|>{\centering\arraybackslash}m{1.8cm}|>{\centering\arraybackslash}m{1cm}|}
\caption{Параметры для настройки ступени ТЗНП\hfill\vspace{-0.5\baselineskip}}\label{tznp:tbl1}\\ 
\hline
\rowcolor{gray!30}
Параметр (Параметр на ИЧМ) & Условное обозначение на схеме & Значение/ Диапазон & Единица измерения & Шаг \\ 
\hline
\endfirsthead
\caption*{\hspace{3pt}\emph{Продолжение таблицы \ref{tznp:tbl1}\hfill\vspace{-0.5\baselineskip}}} \\ % сделано по ГОСТ 2.105 п.6.8.7
\hline
\rowcolor{gray!30}
Параметр (Параметр на ИЧМ) & Условное обозначение на схеме & Значение/ Диапазон & Единица измерения & Шаг \\ 
%\hline
\endhead
%\hline
\endfoot
%\hline
\endlastfoot
%==+t1*ZSPTOC1|SV_LVZSTOC
\centering Ввод функции в работу (Ввод\_функции) & \centering SGF1 & \centering 0 = Не предусмотрено\\1 = Предусмотрено & \centering -- & \centering \arraybackslash -- \\
\hline
\centering Режим оперативного блокирования ЧС (Реж\_БЧС) & \centering SGF8 & \centering 0 = Не предусмотрено\\1 = Предусмотрено & \centering -- & \centering \arraybackslash -- \\
\hline
\centering Режим контроля от БНТ (Реж\_БНТ) & \centering SGF7 & \centering 0 = Не предусмотрено\\1 = Предусмотрено & \centering -- & \centering \arraybackslash -- \\
\hline
\centering Ток срабатывания НП (3I0ср) & \centering 3I0> & \centering 0,10 ... 30,00 & \centering о.е. & \centering \arraybackslash 0,01 \\
\hline
\centering Вывод направленности при АУ (Ненапр\_АУ) & \centering SGF4 & \centering 0 = Не предусмотрено\\1 = Предусмотрено & \centering -- & \centering \arraybackslash -- \\
\hline
\centering Режим направленности (Направленность) & \centering SGF2 & \centering 1 = Ненаправленная\\2 = Прямонаправленная\\3 = Обратнонаправленная & \centering -- & \centering \arraybackslash -- \\
\hline
\centering Режим ОНМ НП при неисправности ЦН (Реж\_БНН\_ОНМ\_НП) & \centering SGF3 & \centering 0 = Блокировка\\1 = Деблокировка & \centering -- & \centering \arraybackslash -- \\
\hline
\centering Режим ПО 3U0 при неисправности ЦН (Реж\_БНН\_3U0) & \centering SGF5 & \centering 0 = Блокировка\\1 = Деблокировка & \centering -- & \centering \arraybackslash -- \\
\hline
\centering Напряжение срабатывания НП (3U0ср) & \centering 3U0> & \centering 2,0 ... 100,0 & \centering В & \centering \arraybackslash 0,1 \\
\hline
\centering Режим контроля от ПО 3U0 (Реж\_3U0) & \centering SGF6 & \centering 0 = Не предусмотрено\\1 = Предусмотрено & \centering -- & \centering \arraybackslash -- \\
\hline
\centering Характеристика срабатывания (Характеристика) & \centering -- & \centering 1 = Кривая F\\2 = Кривая E\\4 = Кривая D\\10 = Кривая B\\11 = Кривая A\\12 = Кривая C\\16 = Независимая\\17 = Кривая RXIDG & \centering -- & \centering \arraybackslash -- \\
\hline
\centering Коэффициент регулировки времени срабатывания зависимой характеристики (К) & \centering -- & \centering 0,05 ... 15,00 & \centering -- & \centering \arraybackslash 0,01 \\
\hline
\centering Независимая выдержка времени срабатывания (Tср) & \centering T1 & \centering 0,00 ... 100,00 & \centering с & \centering \arraybackslash 0,01 \\
\hline
%===t1
\end{longtable}
\normalsize

\vspace{3mm}
\begin{figure}[H]
\centering
\includegraphics[width=0.9\textwidth,height=0.9\textheight,keepaspectratio]{\fbpath/404. ТЗНП СВ/_latex/img/tznpst.pdf}
\captionof{figure}{Функциональная схема ступеней ТЗНП}\label{tznp:img1}
\end{figure}

\end{enumerate}

%%%%%%%%%%%%%%%%%%%%%%%%%%%%%%%%%%%%%%%%%%%%%%%%%%%%%%%%% ОРГАН ОНМ НП %%%%%%%%%%%%%%%%%%%%%%%%%%%%%%%%%%%%%%%%%%%%%%%%

\item Орган направления мощности НП (ОНМ НП)

\begin{enumerate}[label=\arabic{section}.\arabic{subsection}.\arabic{enumi}.\arabic*, labelsep=4pt, leftmargin=0em, itemindent=65pt, parsep=0pt]

\item
В радиальных сетях, когда собственные емкостные токи отдельных присоединений велики и соизмеримы с полным током сети ненаправленная токовая защита нулевой последовательности не может обеспечить селективность действия. Для селективной работы в таких случаях необходимо применение органа направления мощности нулевой последовательности совместно с измерительным токовым органом защиты.

\item
Функция <<ОНМ НП>> автоматически вводится в работу, если введена в работу одна из ступеней ТЗНП и в этой ступени настроен контроль направленности.

\item
В качестве воздействующих величин ОНМ НП использует значения первой гармоники тока нулевой последовательности и напряжения нулевой последовательности.

\item
Направленный элемент срабатывает при выполнении следующих условий:
\begin{itemize}
    \item напряжение НП превышает 1 вольт;
    \item ток НП превышает значение $0,04 \cdot I_{\textsl{ном}}$ НП;
    \item вектор полной мощности нулевой последовательности, определяемый по углу током НП и напряжением НП, находится в области срабатывания (прямое или обратное направление).
\end{itemize}

\item
Для определения области срабатывания задается уставка $\varphi_{\textsl{м.ч.}}$. Отсчет угла максимальной чувствительности производится от вектора напряжения НП и считается положительным, если направление отсчета совпадает с направлением вращения против часовой стрелки. Угол сектора отсчитывается от угла максимальной чувствительности в обе стороны и составляет 85°.

\item 
Сигнал <<ОНМ НП: Прямое>> появляется, если вектор полной мощности НП находится в диапазоне от $\varphi_{\textsl{м.ч.}}$ - 85° до  $\varphi_{\textsl{м.ч.}}$ + 85°. Сигнал <<ОНМ НП: Обратное>> появляется, если вектор полной мощности НП находится в диапазоне от ($\varphi_{\textsl{м.ч.}}$ - 180°) + 85° до ($\varphi_{\textsl{м.ч.}}$ - 180°) - 85°. Если вектор полной мощности НП находится в зоне <<нечувствительности>>, то ни один из вышеуказанных сигналов не появляется.

\item 
Поясняющая диаграмма определения направления мощности приведена на рисунке \ref{tznp:imgonm0}.

\vspace{3mm}
\begin{figure}[H]
\centering
\includegraphics[width=0.6\textwidth,height=0.6\textheight,keepaspectratio]{\fbpath/404. ТЗНП СВ/_latex/img/char.pdf}
\captionof{figure}{Характеристика органа направления мощности нулевой последовательности}\label{tznp:imgonm0}
\end{figure}

\item
Алгоритм работы функции <<ОНМ НП>> выполнен в соответствии с рисунком \ref{tznp:imgonm1}. В таблице \ref{tznp:tblonm} приведены параметры, необходимые для настройки функции <<ОНМ НП>>.

\small
\begin{longtable}{|>{\centering\arraybackslash}m{5.3cm}|>{\centering\arraybackslash}m{3.3cm}|>{\centering\arraybackslash}m{4.2cm}|>{\centering\arraybackslash}m{1.8cm}|>{\centering\arraybackslash}m{1cm}|}
\caption{Параметры для настройки функции <<ОНМ НП>>\hfill\vspace{-0.5\baselineskip}}\label{tznp:tblonm}\\ 
\hline
\rowcolor{gray!30}
Параметр (Параметр на ИЧМ) & Условное обозначение на схеме & Значение/ Диапазон & Единица измерения & Шаг \\ 
\hline
\endfirsthead
\caption*{\hspace{3pt}\emph{Продолжение таблицы \ref{tznp:tblonm}\hfill\vspace{-0.5\baselineskip}}} \\ % сделано по ГОСТ 2.105 п.6.8.7
\hline
\rowcolor{gray!30}
Параметр (Параметр на ИЧМ) & Условное обозначение на схеме & Значение/ Диапазон & Единица измерения & Шаг \\ 
%\hline
\endhead
%\hline
\endfoot
%\hline
\endlastfoot
%==+t1*RDIR1|SV_LVZSTOC
\centering Угол максимальной чувствительности (Фмч) & \centering φмч & \centering -179,0 ... 180,0 & \centering град. & \centering \arraybackslash 0,1 \\
\hline
%===t1
\end{longtable}
\normalsize

\vspace{3mm}
\begin{figure}[H]
\centering
\includegraphics[width=0.9\textwidth,height=0.9\textheight,keepaspectratio]{\fbpath/404. ТЗНП СВ/_latex/img/onmnp.pdf}
\captionof{figure}{Функциональная схема <<ОНМ НП>>}\label{tznp:imgonm1}
\end{figure}

\end{enumerate}

%%%%%%%%%%%%%%%%%%%%%%%%%%%%%%%%%%%%%%%%%%%%%%%%%%%%%%%%% ОРГАН БНТ %%%%%%%%%%%%%%%%%%%%%%%%%%%%%%%%%%%%%%%%%%%%%%%%

\item Блокировка при броске намагничивающего тока НП (БНТ НП)

\begin{enumerate}[label=\arabic{section}.\arabic{subsection}.\arabic{enumi}.\arabic*, labelsep=4pt, leftmargin=0em, itemindent=65pt, parsep=0pt]

\item
Функция блокировки при броске намагничивающего тока НП (<<БНТ НП>>) позволяет заблокировать защиту нулевой последовательности при броске тока намагничивания НП, возникающем при постановке силового трансформатора под напряжение. 
\item
Ввод функции <<БНТ НП>> осуществляется автоматически при активации режим контроля броска тока намагничивания введенной в работу ступени ТЗНП.
\item
Условие пуска выполняется при одновременном наличии следующих условий:
\begin{itemize}
\item отношение действующего значения тока НП второй гармоники (\textsl{3I0 2г}) к действующему значению тока НП первой гармоники (\textsl{3I0}) превышает уставку <<Ih2/Ih1>>;
\item отсутствует срабатывание порогового элемента, реагирующего на превышение действующего значения первой гармоники тока НП (\textsl{3I0}) над уставкой <<Iср>>;
\item сформирован пуск одной из ступеней ТЗНП.
\end{itemize}

\item
Алгоритм работы функции <<БНТ НП>> выполнен в соответствии с рисунком \ref{tznp:img3}. В таблице \ref{tznp:tbl3} приведены параметры, необходимые для настройки функции <<БНТ НП>>.

\vspace{3mm}
\begin{figure}[H]
\centering
\includegraphics[width=0.8\textwidth,height=0.8\textheight,keepaspectratio]{\fbpath/404. ТЗНП СВ/_latex/img/bntnp.pdf}
\captionof{figure}{Функциональная схема <<БНТ НП>>}\label{tznp:img3}
\end{figure}

\small
\begin{longtable}{|>{\centering\arraybackslash}m{5.3cm}|>{\centering\arraybackslash}m{3.3cm}|>{\centering\arraybackslash}m{4.2cm}|>{\centering\arraybackslash}m{1.8cm}|>{\centering\arraybackslash}m{1cm}|}
\caption{Параметры для настройки функции <<БНТ НП>>\hfill\vspace{-0.5\baselineskip}}\label{tznp:tbl3}\\ 
\hline
\rowcolor{gray!30}
Параметр (Параметр на ИЧМ) & Условное обозначение на схеме & Значение/ Диапазон & Единица измерения & Шаг \\ 
\hline
\endfirsthead
\caption*{\hspace{3pt}\emph{Продолжение таблицы \ref{tznp:tbl3}\hfill\vspace{-0.5\baselineskip}}} \\ % сделано по ГОСТ 2.105 п.6.8.7
\hline
\rowcolor{gray!30}
Параметр (Параметр на ИЧМ) & Условное обозначение на схеме & Значение/ Диапазон & Единица измерения & Шаг \\ 
%\hline
\endhead
%\hline
\endfoot
%\hline
\endlastfoot
%==+t1*PHAR1|SV_LVZSTOC
\centering Отношение тока второй гармоники к току основной гармоники (Ih2/Ih1) & \centering I2г/I1г> & \centering 5 ... 100 & \centering \% & \centering \arraybackslash 1 \\
\hline
\centering Ток блокировки НП (3I0ср) & \centering 3I0> & \centering 0,1 ... 15,0 & \centering о.е. & \centering \arraybackslash 0,1 \\
\hline
%===t1
\end{longtable}
\normalsize

\end{enumerate}

%%%%%%%%%%%%%%%%%%%%%%%%%%%%%%%%%%%%%%%%%%%%%%%%%%%%%%%%% ОРГАН АУ %%%%%%%%%%%%%%%%%%%%%%%%%%%%%%%%%%%%%%%%%%%%%%%%

\item Автоматическое ускорение ТЗНП (АУ ТЗНП)

\begin{enumerate}[label=\arabic{section}.\arabic{subsection}.\arabic{enumi}.\arabic*, labelsep=4pt, leftmargin=0em, itemindent=65pt, parsep=0pt]

\item
В устройстве предусмотрен режим автоматического ускорения ступеней ТЗНП. Выбор ускоряемых ступеней защиты осуществляется из соображений надежного охвата, защищаемого энергообъекта.
\item
Предусмотрена возможность автоматического ускорения для любой ступени ТЗНП. Ускоряемая ступень выбирается уставкой SGF1. Автоматическое ускорение срабатывает без выдержки времени.

\item
Алгоритм работы <<АУ ТЗНП>> выполнен в соответствии с рисунком \ref{tznp:imgau}. В таблице \ref{tznp:tblau} приведены параметры, необходимые для настройки функции <<АУ ТЗНП>>.

\begin{figure}[H]
\centering


\includegraphics[width=0.6\textwidth,height=0.6\textheight,keepaspectratio]{\fbpath/404. ТЗНП СВ/_latex/img/au.pdf}
\captionof{figure}{Функциональная схема АУ ТЗНП}\label{tznp:imgau}
\vspace{3mm}
\end{figure}

\small
\begin{longtable}{|>{\centering\arraybackslash}m{5.3cm}|>{\centering\arraybackslash}m{3.3cm}|>{\centering\arraybackslash}m{4.2cm}|>{\centering\arraybackslash}m{1.8cm}|>{\centering\arraybackslash}m{1cm}|}
\caption{Параметры для настройки функции <<АУ ТЗНП>>\hfill\vspace{-0.5\baselineskip}}\label{tznp:tblau}\\ 
\hline
\rowcolor{gray!30}
Параметр (Параметр на ИЧМ) & Условное обозначение на схеме & Значение/ Диапазон & Единица измерения & Шаг \\ 
\hline
\endfirsthead
\caption*{\hspace{3pt}\emph{Продолжение таблицы \ref{tznp:tblau}\hfill\vspace{-0.5\baselineskip}}} \\ % сделано по ГОСТ 2.105 п.6.8.7
\hline
\rowcolor{gray!30}
Параметр (Параметр на ИЧМ) & Условное обозначение на схеме & Значение/ Диапазон & Единица измерения & Шаг \\ 
%\hline
\endhead
%\hline
\endfoot
%\hline
\endlastfoot
%==+t1*AUAPSSU1|SV_LVZSTOC
\centering Выбор ускоряемой ступени (Выбор\_ст) & \centering SGF1 & \centering 1 = 1 ступень\\2 = 2 ступень\\3 = 3 ступень\\4 = 4 ступень\\5 = 1 или 3 ступень\\6 = 2 или 4 ступень & \centering -- & \centering \arraybackslash -- \\
\hline
%===t1
\end{longtable}
\normalsize

\end{enumerate}
%%%%%%%%%%%%%%%%%%%%%%%%%%%%%%%%%%%%%%%%%%%%%%%%%%%%%%%%% КОНЕЦ ОРГАН АУ %%%%%%%%%%%%%%%%%%%%%%%%%%%%%%%%%%%%%%%%%%%%%%%%

%%%%%%%%%%%%%%%%%%%%%%%%%%%%%%%%%%%%%%%%%%%%%%%%%%%%%%%%% ОРГАН ОУ %%%%%%%%%%%%%%%%%%%%%%%%%%%%%%%%%%%%%%%%%%%%%%%%

\item Оперативное ускорение ТЗНП (ОУ ТЗНП)

\begin{enumerate}[label=\arabic{section}.\arabic{subsection}.\arabic{enumi}.\arabic*, labelsep=4pt, leftmargin=0em, itemindent=65pt, parsep=0pt]

\item
Оперативное ускорение предусматривает ускорение любой ступени ТЗНП. Оперативное ускорение может вводиться:
\begin{itemize}
    \item при выводе основной защиты оборудования;
    \item при опробовании оборудования рабочим напряжением;
    \item при операциях с разъединителями на противоположном конце ВЛ.
\end{itemize}

\item
В данном алгоритме при введенном ключе <<Ввод ОУ ТЗНП>> предусмотрена возможность оперативного ускорения любой ступени ТЗНП. Выбор ускоряемой ступени осуществляется программным переключателем SGF1 <<Выбор\_ст\_ОУ>>. При введенной функции и появлению сигнала пуска выбранной ступени по истечению выдержки времени, определяемой уставкой T1, формируется сигнал срабатывания.

\item
Алгоритм работы <<ОУ ТЗНП>> выполнен в соответствии с рисунком \ref{tznp:imgou}. В таблице \ref{tznp:tblou} приведены параметры, необходимые для настройки функции <<ОУ ТЗНП>>.

\begin{figure}[H]
\centering

\includegraphics[width=0.8\textwidth,height=0.8\textheight,keepaspectratio]{\fbpath/404. ТЗНП СВ/_latex/img/ou.pdf}
\captionof{figure}{Функциональная схема ОУ ТЗНП}\label{tznp:imgou}
\vspace{3mm}
\end{figure}

\small
\begin{longtable}{|>{\centering\arraybackslash}m{5.3cm}|>{\centering\arraybackslash}m{3.3cm}|>{\centering\arraybackslash}m{4.2cm}|>{\centering\arraybackslash}m{1.8cm}|>{\centering\arraybackslash}m{1cm}|}
\caption{Параметры для настройки функции <<ОУ ТЗНП>>\hfill\vspace{-0.5\baselineskip}}\label{tznp:tblou}\\ 
\hline
\rowcolor{gray!30}
Параметр (Параметр на ИЧМ) & Условное обозначение на схеме & Значение/ Диапазон & Единица измерения & Шаг \\ 
\hline
\endfirsthead
\caption*{\hspace{3pt}\emph{Продолжение таблицы \ref{tznp:tblou}\hfill\vspace{-0.5\baselineskip}}} \\ % сделано по ГОСТ 2.105 п.6.8.7
\hline
\rowcolor{gray!30}
Параметр (Параметр на ИЧМ) & Условное обозначение на схеме & Значение/ Диапазон & Единица измерения & Шаг \\ 
%\hline
\endhead
%\hline
\endfoot
%\hline
\endlastfoot
%==+t1*OPAPSSU1|SV_LVZSTOC
\centering Выбор оперативно ускоряемой ступени (Выбор\_ст\_ОУ) & \centering SGF1 & \centering 1 = 1 ступень\\2 = 2 ступень\\3 = 3 ступень\\4 = 4 ступень\\5 = 1 или 3 ступень\\6 = 2 или 4 ступень & \centering -- & \centering \arraybackslash -- \\
\hline
\centering Выдержка времени срабатывания оперативного ускорения (T\_ОУ) & \centering T1 & \centering 0,00 ... 3,00 & \centering с & \centering \arraybackslash 0,01 \\
\hline
%===t1
\end{longtable}
\normalsize

\end{enumerate}
%%%%%%%%%%%%%%%%%%%%%%%%%%%%%%%%%%%%%%%%%%%%%%%%%%%%%%%%% КОНЕЦ ОРГАН ОУ %%%%%%%%%%%%%%%%%%%%%%%%%%%%%%%%%%%%%%%%%%%%%%%%

\end{enumerate}
